% Chapter 12 -- Disks (Part II, Reference).  Sources (Law 1):
% os/milton/kmain.c (flair_tenant_mount_volume: the data volume mounted as
% drive A; desktop_db_volume_label: the disk icon's name is the volume label
% read off the disk at start-up, an unlabelled volume shows a bare band;
% FLAIR-FAT-MOUNT-FAIL: no data disk, no Finder windows, the clock unset;
% TRASH and DESKTOP.DB made when missing); os/flair/finder_menu.c (Eject,
% Erase Disk, Get Info always dimmed; Open lit by a selected folder);
% os/flair/finder_ops.c (a drop on the disk icon moves to the top level);
% spec/hardware.json (the drives); the data volume as minted by the
% Makefile (FLAIR_FIXTURES: a 1.44 MB FAT12 diskette image, label INITECH).
% The console figures are captures of the InitechDOS prompt reached with
% Exit to InitechDOS (APPS\LABEL, APPS\TREE, APPS\ATTRIB; CHKDSK in
% Chapter 18).  Model: Macintosh User's Guide 1991 Ch. 9 "Hard Disks and
% Floppy Disks" (apple-macintosh.md Sec 8.2).
% GATES: test-flair-data-volume test-flair-desktop-icons test-desktop-db test-finder-menu test-flair-exit-to-dos test-label-vol test-xcopy-tree test-attrib test-chkdsk test-pubs-flair-screens
\Gates{test-flair-data-volume, test-flair-desktop-icons, test-desktop-db,
  test-finder-menu, test-flair-exit-to-dos, test-label-vol, test-xcopy-tree,
  test-attrib, test-chkdsk, test-pubs-flair-screens}
\chapter{Disks}
\label{ch:disks}
\index{disks}

Your work is kept on disks.  This chapter describes the two disks
InitechOS uses, how FLAIR shows them, how much they hold, and how to look
after the files on them.

\section{The startup disk and the data disk}
\index{startup disk}\index{data disk}
InitechOS uses two disks, each in its own drive.

\SideHead{Startup disk}Holds InitechDOS and FLAIR.  The computer loads
them from it when you turn it on.  FLAIR does not show the startup disk on
the desktop, and you do not keep your own files on it.

\SideHead{Data disk}Holds your documents, the applications that come with
FLAIR, and the InitechDOS programs that run at the command prompt.  FLAIR
shows it on the desktop as the \texttt{INITECH} icon, and InitechDOS calls
it \term{drive A}.

Both disks must stay in their drives while FLAIR is running.  Appendix~A
lists every file on the data disk as it is delivered.

\section{The disk icon}
\index{disk icon}\index{INITECH icon}
The data disk is the icon at the top right of the desktop.
\begin{itemize}
\item Double-click it to open the \texttt{Drive A Files} window, which
  shows the top level of the disk.
\item Drag a file or folder onto it to move the item to the top level of
  the disk.
\item Drag the icon itself to move it on the desktop.  FLAIR remembers its
  place.
\end{itemize}
The name under the icon is the disk's \term{volume label}, the name stored
on the disk itself.  The data disk is delivered with the label
\texttt{INITECH}.

\section{What you cannot do with the disk icon}
\index{Eject command}\index{Erase Disk command}\index{Get Info command}
Some commands that act on disks are in the Finder's menus but are dimmed in
this version of FLAIR:
\begin{itemize}
\item \menuitem{Eject} (\menuname{File} menu): the data disk cannot be
  ejected while FLAIR runs.  Shut the computer down before you take a disk
  out of its drive.
\item \menuitem{Erase Disk} (\menuname{Special} menu): FLAIR does not
  erase disks.
\item \menuitem{Get Info} (\menuname{File} menu): FLAIR does not show
  information about a disk, folder or file.  To see how much space the disk
  has, use \file{CHKDSK} at the InitechDOS prompt (see ``How much a disk
  holds,'' below).
\end{itemize}
Selecting the disk icon does not make \menuitem{Open} available in the
\menuname{File} menu; double-click the icon to open it.  \menuitem{Open} is
available when a folder is selected in a Finder window.

\FlairScreen[scale=4, label=fig:kfilesel]{k_filesel}{The File menu with
the APPS folder selected: Open and Close Window are available}

\section{How much a disk holds}
\index{disk space}\index{CHKDSK}
The data disk is a 1.44~MB diskette: it holds 1,457,664 bytes, about 1,424
kilobytes.  As delivered, about 1,220 kilobytes of it are free.

FLAIR does not show the free space on a disk.  To see it, leave FLAIR and
run \file{CHKDSK} from the \file{APPS} directory; Figure~\ref{fig:doschk} in
Chapter~\ref{ch:dos} shows its report.  The \texttt{bytes available on disk}
line is the free space.

Remember that what you put in the Trash still takes up space on the disk
until you empty the Trash.  The alert that asks you to confirm
\menuitem{Empty Trash} tells you how much space the items in the Trash use.

When the disk is full, a new folder cannot be made and a document cannot
be saved.  Empty the Trash, and move documents you no longer need onto
another disk at the InitechDOS prompt.

\section{The name of the disk}
\index{volume label}\index{LABEL}
To see or change the volume label, use the \file{LABEL} program at the
InitechDOS prompt.
\begin{TaskSteps}
\Step Choose \menuitem{Exit to InitechDOS} from the \menuname{Special} menu,
click \texttt{Exit}, and press \key{Enter} twice to keep the date and the
time.
\Step Type \type{APPS\textbackslash LABEL} and press \key{Enter}.
\Result{\file{LABEL} shows the label of drive A and asks for a new one.}
\end{TaskSteps}

\FlairText[fig:klabel]{k_label}{LABEL, showing the name of the data disk}

\begin{TaskSteps}[resume]
\Step Type a new name of up to 11 characters and press \key{Enter}, or
press \key{Enter} alone to leave the label as it is.
\Result{The \doctitle{InitechDOS Reference} describes the questions
\file{LABEL} asks.}
\end{TaskSteps}
FLAIR reads the label when it starts.  A new label appears under the disk
icon the next time you start FLAIR.  A disk with no label is shown with an
empty name band under its icon.

\section{The folders of the disk}
\index{TREE}\index{folders!listing}
The \file{TREE} program lists every folder on the disk, which InitechDOS
calls directories.  Type \type{APPS\textbackslash TREE} at the prompt.

\FlairText[fig:ktree]{k_tree}{The folders of the data disk, listed by TREE}

The list shows the \file{TRASH} directory, which FLAIR uses to hold the
Trash and does not show in its windows (Chapter~\ref{ch:finder}).

\section{Locking a file}
\index{locked items}\index{ATTRIB}\index{read-only files}
A locked file cannot be changed, thrown away or replaced.  Lock a file you
want to keep safe from mistakes, such as a finished report.  You lock and
unlock files at the InitechDOS prompt with the \file{ATTRIB} program;
InitechDOS calls a locked file \textit{read-only}.
\begin{TaskSteps}
\Step At the InitechDOS prompt, type this line and press \key{Enter}:\newline
\type{APPS\textbackslash ATTRIB +R README.TXT}
\Result{\file{README.TXT} is locked.}
\Step Type this line and press \key{Enter}:\newline
\type{APPS\textbackslash ATTRIB README.TXT}
\Result{\file{ATTRIB} shows the file's attributes.  The \texttt{R} at the
left means read-only: the file is locked.}
\end{TaskSteps}

\FlairText[fig:kattrib]{k_attrib}{README.TXT locked with ATTRIB}

To unlock the file, type \type{APPS\textbackslash ATTRIB -R README.TXT}.

In FLAIR, a locked file is shown like any other file.  When you drag it to
the Trash it slides back, and an alert tells you that the item is locked
(Figure~\ref{fig:locked}).  An application cannot save a document under the
name of a locked file: InitechText reports that the document could not be
saved (Figure~\ref{fig:nosave}).  You can still open a locked document and
read it, and you can move a locked file from one folder to another.

\section{If FLAIR starts without the data disk}
\index{data disk!missing}
If the data disk is not in its drive when the computer starts, FLAIR still
starts, but it cannot show any files.

\FlairScreen[scale=4, label=fig:knodisk]{k_nodisk}{FLAIR started without
its data disk: the disk icon has no name, and the clock reads 12:00 am}
\Operation{FLAIR-FAT-MOUNT-FAIL}

The disk icon has no name, and the clock shows 12:00 am, because FLAIR
reads the time through InitechDOS, which has no disk to work with.  Shut the
computer down, put the data disk in its drive, and start the computer
again.

\section{Making copies of your work}
\index{copying files}\index{backup copies}
FLAIR moves files but does not copy them.  To keep a copy of your
documents on another disk, use the \file{COPY} command or the \file{XCOPY}
program at the InitechDOS prompt; the \doctitle{InitechDOS User's Guide}
shows you how.  Make the copies at the end of each day's work, and keep
them away from the computer.
